%!TEX root=../../note.tex

\subsection{Convergent Sequences}
\subsubsection{The Limit of a Sequence}

A sequence of real numbers is a function $X:\N\to\R$. We write $x_n=X(n)$ for its $n$th term and denote the sequence by $\set{x_n}_{n=1}^{\infty}$.

\begin{definition}
    A sequence $\set{x_n}$ in $\R$ converges to $a\in\R$ if, for every $\epsilon>0$, there is an index $n_\epsilon\in\N$ such that
    \[
        n\geq n_\epsilon \quad\Longrightarrow\quad \abs{x_n-a}<\epsilon.
    \]
    We then write $\lim_{n\to\infty}x_n=a$.
\end{definition}

In symbols,
\[
    x_n\to a \iff (\forall\epsilon>0)(\exists n_\epsilon\in\N)(\forall n\geq n_\epsilon)\,\bigl[\abs{x_n-a}<\epsilon\bigr].
\]
In words: for every tolerance $\epsilon>0$, all sufficiently late terms lie in $(a-\epsilon,a+\epsilon)$. The index $n_\epsilon$ may depend on $\epsilon$, and once $n\geq n_\epsilon$ the inequality must hold for \emph{every} later term, even if earlier terms lie outside the interval.

The points below show $x_n=1/n$. For the illustrated tolerance $\epsilon=0.3$, every term from $n_\epsilon=4$ onward lies in the shaded band around $a=0$.

\begin{center}
\begin{tikzpicture}[x=0.55cm,y=2.4cm,>=Stealth]
    \fill[red!8] (0,0) rectangle (11,0.3);
    \draw[densely dashed,red!75!black] (0,0.3) -- (11,0.3) node[right] {$a+\epsilon$};
    \draw[->] (0,0) -- (11.5,0) node[right] {$n$};
    \draw[->] (0,0) -- (0,1.15) node[above] {$x_n$};
    \node[left] at (0,0) {$a=0$};
    \draw[densely dashed,blue!70!black] (4,0) -- (4,0.95);
    \node[below] at (4,0) {$n_\epsilon=4$};
    \foreach \n in {1,...,10} {
        \fill[teal!70!black] (\n,{1/\n}) circle (1.7pt);
    }
    \node[teal!70!black,anchor=west] at (1.3,0.85) {$x_n=1/n$};
\end{tikzpicture}
\end{center}

For $a\in\R$, the $\epsilon$-neighbourhood of $a$ is
\[
    V_\epsilon(a)=\set{x\in\R:\abs{x-a}<\epsilon}=(a-\epsilon,a+\epsilon).
\]
Thus $x_n\to a$ precisely when, for every $\epsilon>0$, there is $n_\epsilon\in\N$ such that $x_n\in V_\epsilon(a)$ for all $n\geq n_\epsilon$.

\begin{example}
    Prove that $\lim_{n\to\infty}\frac1n=0$.
\end{example}
\mdfsetup{nobreak=true}
\begin{proof}
    Let $\epsilon>0$. By the Archimedean Property, choose $n_\epsilon\in\N$ with $1/n_\epsilon<\epsilon$. For every $n\geq n_\epsilon$,
    \[
        \abs{\frac1n-0}=\frac1n\leq\frac1{n_\epsilon}<\epsilon.
    \]
    Hence $1/n\to0$.
\end{proof}
\mdfsetup{nobreak=false}

\begin{example}
    Prove that $\lim_{n\to\infty}\frac{4n+3}{n+1}=4$.
\end{example}
\begin{scratch}
    Simplify the distance first, then see how large $n$ must be:
    \[
        \abs{\frac{4n+3}{n+1}-4}=\abs{\frac{-1}{n+1}}=\frac1{n+1}<\frac1n.
    \]
    So it is enough that $\frac1n<\epsilon$. By Archimedes pick $n_\epsilon$ with $\frac1{n_\epsilon}<\epsilon$; every $n\geq n_\epsilon$ then works.
\end{scratch}
\begin{proof}
    Let $\epsilon>0$. By the Archimedean Property, choose $n_\epsilon\in\N$ with $1/n_\epsilon<\epsilon$. For every $n\geq n_\epsilon$,
    \[
        \abs{\frac{4n+3}{n+1}-4}
        =\frac1{n+1}\leq\frac1n\leq\frac1{n_\epsilon}<\epsilon.
    \]
    Therefore the sequence converges to $4$.
\end{proof}

\subsubsection{Uniqueness of the Limit}
\begin{proposition}
    If $x,y\in\R$ and $\abs{x-y}<\epsilon$ for every $\epsilon>0$, then $x=y$.
\end{proposition}
\begin{proof}
    Suppose $x\ne y$. Then $\epsilon=\abs{x-y}/2>0$, but $\abs{x-y}>\epsilon$, contradicting the hypothesis.
\end{proof}

\begin{theorem}
    A sequence in $\R$ can have at most one limit.
\end{theorem}
\begin{scratch}
    By the preceding proposition, it suffices to show $\abs{a-b}<\epsilon$ for every $\epsilon>0$. The only link between $a$ and $b$ is the sequence, so pass through a term $x_n$:
    \[
        \abs{a-b}\leq\abs{a-x_n}+\abs{x_n-b}.
    \]
    To make the sum $<\epsilon$, make each piece $<\epsilon/2$: that needs $n\geq n_1$ for the first and $n\geq n_2$ for the second, so take $n\geq\max\set{n_1,n_2}$.
\end{scratch}
\begin{proof}
    Suppose $x_n\to a$ and $x_n\to b$. Given $\epsilon>0$, choose $n_1,n_2\in\N$ so that $\abs{x_n-a}<\epsilon/2$ whenever $n\geq n_1$ and $\abs{x_n-b}<\epsilon/2$ whenever $n\geq n_2$. For $n\geq\max\set{n_1,n_2}$, the triangle inequality gives
    \[
        \abs{a-b}\leq\abs{a-x_n}+\abs{x_n-b}<\frac\epsilon2+\frac\epsilon2=\epsilon.
    \]
    This holds for every $\epsilon>0$, so the preceding proposition gives $a=b$.
\end{proof}

\subsubsection{Divergence and Boundedness}
A sequence diverges if it does not converge to any real number. For example, consider $\set{x_n}=\set{0,1,0,1,\ldots}$, with $x_n=0$ for odd $n$ and $x_n=1$ for even $n$.

Both $0$ and $1$ keep appearing arbitrarily late, so the terms cannot all remain close to one proposed limit.

Suppose $x_n\to a$ and take $\epsilon=1/2$. There is an $n_\epsilon$ such that $\abs{x_n-a}<1/2$ for every $n\geq n_\epsilon$. Choose an odd and an even index beyond $n_\epsilon$. The odd term gives $\abs{a}<1/2$, hence $-1/2<a<1/2$. The even term gives $\abs{1-a}<1/2$, hence $1/2<a<3/2$. These intervals do not overlap, a contradiction. Thus this sequence diverges.

\begin{definition}
    A sequence $\set{x_n}$ is bounded if there is an $M\in\R$ such that $\abs{x_n}\leq M$ for every $n\in\N$.
\end{definition}

\begin{theorem}
    Every convergent sequence of real numbers is bounded.
\end{theorem}
\begin{idea}
    The tail stays within $1$ of the limit; only finitely many terms come before it, and a finite set has a maximum.
\end{idea}
\begin{proof}
    Suppose $x_n\to a$. Taking $\epsilon=1$, choose $n_1\in\N$ so that $\abs{x_n-a}<1$ for all $n\geq n_1$. For these terms, the triangle inequality gives
    \[
        \abs{x_n}\leq\abs{x_n-a}+\abs{a}<1+\abs{a}.
    \]
    The remaining terms form a finite set. Thus
    \[
        M=\max\set{\abs{x_1},\ldots,\abs{x_{n_1-1}},1+\abs{a}}
    \]
    bounds every term of the sequence (when $n_1=1$, take $M=1+\abs{a}$).
\end{proof}

The converse fails: $\set{0,1,0,1,\ldots}$ is bounded but, as shown above, diverges.


\begin{example}
    Consider the following sequences. 
    \begin{itemize}
        \item $\set{a_n } = \set{\frac{n^2 + 1 }{n }}$
        
        Take a look at this sequence: 
        \begin{equation*}
            \abs{\frac{n^2 + 1 }{n }} = \abs{n + \frac{1}{n }} = n + \frac{1}{n}
        \end{equation*}
    
        For any $M > 0$, there exist $n_m \in \N$ such that $n_M > M$.  Thus, 
        \begin{equation*}
            \abs{a_{n_M}} > M. 
        \end{equation*}

        \item $\set{a_n } = \set{\paren{-1}^n n}$
        
        Take a look at this sequence:

        \begin{equation*}
            \abs{a_n } = \abs{\paren{-1}^n n } = \abs{\paren{-1}^n }\abs{n } = \abs{n } > M
        \end{equation*}
    \end{itemize}
\end{example}

\subsubsection{Divergence to $\pm\infty$}
\begin{definition}
    A sequence $\set{x_n }$ diverges to $+\infty$, that is, $\lim_{x_n } = \infty$, if for all real number $M > 0$, there exists $n_M \in \N$
    such that $a_n > M$ for all $n \geq n_M$. 
\end{definition}

\begin{definition}
    A sequence $\set{x_n }$ diverges to $-\infty$, that is, $\lim_{x_n } = -\infty$, if for all real number $M > 0$, there exists $n_M \in \N$
    such that $a_n < - M$ for all $n \geq n_M$. 
\end{definition}
We write $\lim_{n\to\infty}x_n=\pm\infty$, but such a sequence is still divergent: it does not converge to any real number.


\begin{example}
    Prove that $\lim_{n \to \infty } \frac{n^2}{1 - n} = -\infty$. 
\end{example}
\begin{proof}
    Let $M > 0$. 
\end{proof}

\begin{example}
    Prove that $\lim_{x \to \infty} \sqrt{n } + 1 = + \infty$. 
\end{example}
\begin{proof}
    The goal is to let $\sqrt{n } > m - 1$
\end{proof}

\begin{theorem}
    If $\lim_{n \to \infty} x_n = a$ and $\lim_{n \to \infty} y_n = b$, then 
    \begin{enumerate}
        \item $(\forall c \in \R)\,[\lim_{n \to \infty} c = c]$
        \item $(\forall c \in \R)\,[\lim_{n \to \infty} cx_n = ca]$
    \end{enumerate}
\end{theorem}